\documentclass[12pt,a4paper]{article}

\usepackage{amsmath,amssymb}
\usepackage{graphicx}
\usepackage{geometry}
\geometry{left=3cm,right=3cm,top=2.5cm,bottom=2.5cm}
\usepackage{setspace}
\onehalfspacing
\usepackage{parskip}
\setlength{\parskip}{0.7em}
\usepackage[T1]{fontenc}
\usepackage{lmodern}
\usepackage{microtype}
\usepackage{placeins}
\usepackage{hyperref}
\hypersetup{colorlinks=true, linkcolor=black, citecolor=black, urlcolor=black}

\title{\LARGE\textbf{The Climate Standoff}}
\author{}
\date{}

\begin{document}

\maketitle
\thispagestyle{empty}
\newpage

% ============================================================
\section*{Introduction}
\addcontentsline{toc}{section}{Introduction}
% ============================================================

Threshold public goods models explain coordination failure by showing that collective benefits depend on reaching a critical mass, while the presence of multiple equilibria makes that mass difficult to attain. However, the most significant shifts in global climate policy suggest that enforcement pressure and technological progress are endogenous byproducts of domestic adoption rather than formal prerequisites. The political viability of the European Union's Carbon Border Adjustment Mechanism (CBAM) was not a result of a global treaty. Instead, it was made possible only after domestic decarbonization created the incentive to export climate standards through trade. Similarly, the United States' Inflation Reduction Act (IRA) reshaped global investment through massive subsidies, forcing a competitive policy response across Europe and Asia. In both cases, the incentives to participate were generated by the strategic decisions of influential actors moving first, which fundamentally altered the payoff matrix for everyone else.

The strategic problem is best understood as a multilateral game of chicken with highly asymmetric players. In this environment, the blocs are not driving identical cars toward the same obstacle. A small number of players are so systemically important that their failure to turn carries disproportionate global consequences, effectively forcing them into a high stakes standoff. Others face weaker brakes because short political horizons and high transition costs make early adjustment more difficult. Some are closer to the cliff because climate damages arrive faster or hit more vulnerable sectors, while others can afford to delay longer despite their pivotal influence over the collective outcome.

The coordination problem is therefore shaped not only by whether blocs eventually swerve, but by who can afford to wait, who bears the greatest crash risk, and whose movement most strongly changes the incentives of everyone else. Standard agreement based models abstract from these dynamics and consequently struggle to explain why cooperation can deepen even when formal treaties stagnate. By treating pressure, cost declines, and threshold progress as endogenous strategic outcomes, we can account for the timing, escalation, and asymmetry that define modern climate geopolitics.

This paper develops a finite horizon dynamic Markov game to analyze these interactions. In this framework, heterogeneous blocs face binary adoption decisions where transition costs decline through learning and political pressure rises endogenously with each new adopter. Collective benefits are only realized once a weighted participation threshold is met. By modeling these variables as results of strategic interaction, this framework captures how the early moves of anchor players can trigger a tipping point, transforming a dangerous game of delay into a self reinforcing race toward cooperation.

% ============================================================
\section{Literature Review}
% ============================================================

\subsection{Equilibrium Concepts: Nash to Quantal Response}

Standard game theoretic analyses of climate cooperation typically rely on Nash equilibrium or its refinements. Barrett (1994) employs both stage game Nash equilibrium and renegotiation proof repeated game equilibrium to characterize self enforcing international environmental agreements. Harstad (2016) works within a dynamic framework, selecting Markov perfect equilibrium because strategies depending only on payoff relevant state variables are both analytically tractable and behaviorally plausible. These solution concepts share a common feature in that players are assumed to best respond exactly to the strategies of others without systematic noise or payoff sensitive error.

Within the specific class of models that Barrett inaugurates, this assumption generates a well known fragility. Because participation decisions are discrete and stable coalitions are determined by internal and external stability conditions, small changes in parameters can induce large, discontinuous shifts in equilibrium coalition size. Outcomes move abruptly between near full cooperation and near complete defection, which limits the usefulness of such frameworks for analyzing how coordination responds to marginal changes in underlying incentives.

The empirical design of this dissertation requires equilibrium behavior that varies smoothly with parameters. Quantal response equilibrium provides this structure by assuming that players choose actions probabilistically. Choice probabilities increase in expected payoffs while remaining stochastic, governed by a rationality parameter lambda. As lambda increases, behavior approaches the deterministic best responses of Nash equilibrium. As lambda decreases, choices converge toward randomness. This formulation ensures that small changes in payoffs lead to small changes in behavior. For example, a marginal increase in adoption costs reduces the probability of adoption incrementally rather than triggering a discrete shift from full participation to none.

The specific value of lambda used here is calibrated to the model, and the fixed point algorithm converges to a consistent solution across initializations. While the necessary trade off is the calibration of lambda itself, the parameter carries a clear substantive interpretation. It captures the noise inherent in international negotiations where leaders are constrained by domestic win sets. In this framework, stochastic choice is not merely a technical device to ensure smoothness. Instead, it reflects a world of incomplete information, bureaucratic frictions, and internal political pressures that prevent governments from acting as perfectly rational best responders. By situating the model within this two level strategic context, the approach provides a more realistic representation of how sovereign states navigate the friction between international commitments and domestic survival.

However, while quantal response resolves the problem of equilibrium brittleness, it does not by itself explain how marginal shifts in incentives translate into the self reinforcing tipping dynamics observed in real world climate transitions. This raises a further question regarding when smooth changes in strategic incentives accumulate into cascading adoption rather than stalling at partial cooperation.

\subsection{Strategic Tipping and Diffusion Dynamics}

While the previous section resolves how equilibrium behavior should respond smoothly to changing incentives, an additional question remains regarding the nature of global transitions. Recent work on tipping points and technology diffusion in the post Paris literature has increasingly focused on nonlinear dynamics in climate cooperation. Farmer et al.\ (2019) introduce the concept of sensitive intervention points in socioeconomic systems, arguing that small and well targeted interventions can exploit amplification mechanisms to trigger cascading low carbon transitions. Otto et al.\ (2020) identify specific social subsystems where disruptive change driven by positive feedback could rapidly reduce global emissions. Furthermore, Sharpe and Lenton (2021) extend this logic into a framework of upward scaling tipping cascades, in which activating one socio-technical tipping point increases the likelihood of triggering another at a larger scale. Way et al.\ (2022) incorporate endogenous cost declines to demonstrate how fast learning curves for renewables can significantly alter mitigation pathways and deliver large net savings.

These contributions capture tipping behavior observed in real world energy transitions, but they share a structural limitation. Adoption decisions are typically embedded within aggregate cost functions or representative agent frameworks rather than arising from strategic interaction between heterogeneous countries. While existing research correctly identifies the existence of reinforcing policy loops, the absence of a game theoretic foundation makes it difficult to disentangle the underlying drivers. It remains unclear whether a cascade is driven by falling costs, shifting expectations, or peer pressure from early adopters. This lack of strategic structure limits the ability of existing models to explain why some countries adopt early while others delay, or to identify precisely when gradual progress transitions into self reinforcing coordination rather than stalling.

This ambiguity is compounded by the fact that tipping behavior in formally specified models often depends heavily on chosen parameter values. This sensitivity makes it difficult to separate structurally robust dynamics from outcomes that are merely calibration specific. Taken together, these limitations point to two largely unresolved problems in the literature. These include the absence of a strategic framework that endogenizes tipping through asymmetric interaction and the lack of a disciplined way to distinguish structurally robust findings from results that depend on specific model calibrations.

\subsection{Gaps and Contribution}

Two gaps run through the literature reviewed above.

First, few existing frameworks model the interaction between declining costs, endogenous political pressure, and asymmetric influence as joint determinants of climate cooperation. Agreement-based models typically treat enforcement as exogenous, while diffusion models embed adoption in aggregate cost functions rather than strategic interaction. As a result, they struggle to explain how early adoption by influential players reshapes the incentives of strategic holdouts over time.

Second, while game-theoretic models of climate cooperation are known to be sensitive to parameterisation, few studies distinguish data-identified cross-bloc structure from the latent cardinal strength of the underlying payoff channels. Most computational approaches either calibrate to a single benchmark or rely on ad hoc sensitivity checks, making it difficult to separate structurally robust findings from calibration-driven ones.

This paper addresses both limitations. It develops a finite-horizon dynamic Markov game in which heterogeneous blocs make irreversible adoption decisions under declining transition costs, rising endogenous pressure, and threshold-based stabilisation benefits. Behaviour is characterised using quantal response equilibrium, allowing smooth comparative statics and realistic strategic delay.

Methodologically, the contribution lies in separating calibration into ordinal structure and cardinal channel strength. The ordinal layer is pinned down using macroeconomic and emissions data, preserving the empirical asymmetries across blocs. The cardinal layer is then estimated using a simulated method of moments procedure that matches stylised features of the contemporary climate transition, including EU leadership, strategic free-riding, and delayed tipping.

This produces a single fitted strategic baseline from which the analysis identifies which policy levers most effectively increase the likelihood of earlier and more reliable coordination. By combining local sweeps with global sensitivity analysis, the framework distinguishes locally powerful channels from structurally robust levers that remain influential once the wider strategic environment is allowed to move.

% ============================================================
\section{The Model}
% ============================================================

The model is a finite-horizon dynamic Markov game in which countries repeatedly decide whether to irreversibly adopt green technology. The strategic problem is forward-looking: current decisions affect future incentives by changing the aggregate adoption state. Because adoption decisions are simultaneous and unobserved within periods, players face strategic uncertainty over future adoption paths. Equilibrium behaviour is therefore characterised using a dynamic quantal response equilibrium, which allows probabilistic choice and ensures well-defined comparative statics.

\subsection{Model Primitives}

This section outlines the core building blocks of the model: the players and the game structure.

\subsubsection{Players}

\[
  i \in \{\text{US},\, \text{EU},\, \text{CN},\, \text{RoW}\}
\]

Each player represents a geopolitical bloc with distinct economic structures, emissions profiles, and strategic incentives. Using representative blocs rather than individual countries is standard in the climate cooperation literature, as it preserves the key asymmetries that shape global outcomes while maintaining tractability.

The United States, Europe, and China together account for a huge chunk of global emissions, renewable manufacturing capacity, and dominate geopolitical influence. The Rest of the World aggregates developing and emerging economies whose individual influence is limited but whose collective climate vulnerability is high.

This structure captures a central asymmetry of climate cooperation: players with the greatest capacity to influence global outcomes often face weaker direct climate incentives, while those with the strongest incentives face limited capacity to drive coordination.

\subsubsection{Time and Actions}

Time is discrete and finite, indexed by:
\[
  t = 1, 2, \ldots, T
\]
The finite horizon captures the urgency of climate action, reflecting deadlines beyond which damages become irreversible or catastrophically costly.

In each period $t$, players simultaneously choose an adoption action:
\[
  G_{it} \in \{0, 1\}
\]
where $G_{it} = 1$ denotes adoption of green technology and $G_{it} = 0$ denotes continued delay. Adoption is irreversible:
\[
  G_{it} \leq G_{i,t+1}
\]
Once a player adopts, it remains an adopter in all subsequent periods. This reflects the sunk and long-lived nature of energy infrastructure investments and policy commitments.

\[
  \tau_i = \min\{t : G_{it} = 1\}, \qquad \tau_i = \infty \;\text{if never}
\]

which is the object of strategic choice. What time is optimal to adopt? If ever?

But crucially not all players' choices affect the state in the same way. This is represented by players being assigned an influence weight.

\subsubsection{Player Influence Weights}

Each player $i$ is assigned an influence weight $w_i > 0$, representing the share of global impact it contributes when adopting green technology. These weights reflect differences in economic scale, energy consumption, and geopolitical leverage. Major emitters receive relatively large weights, while the aggregated RoW bloc receives a smaller but meaningful share. Weights are normalised such that:
\[
  \sum_i w_i = 1
\]
allowing weighted adoption to be interpreted directly as the proportion of global influence committed to the transition. Influence asymmetry is crucial in determining how easily the threshold can be reached. A transition in which only small players adopt may fail to cross the threshold, whereas early adoption by the US or China can trigger cascades by substantially increasing weighted adoption in a single period.

\subsubsection{State: Defining the Game State}

\[
  s_t = (t,\, W)
\]

\textit{Implementation note: While strategies are written as functions of $(t, W)$ for notational clarity, the algorithm tracks the full adoption vector $G$ to correctly compute transition probabilities. At each state, only players who have not yet adopted contribute to strategic uncertainty in the next period. This distinction matters when multiple adoption configurations yield the same aggregate $W$.}

$W$ is the cumulative weighted share of adoption:
\[
  W \equiv \sum_i w_i G_i
\]

This simply defines the state of the game: it captures time (which round of the game are players in?) and the cumulative adoption (what share of the world has committed to green energy?).

At the start of each round, players look at the current state $W_t$ to decide their next move. Because everyone acts at the same time, no one knows what the others are doing until the round is over. Once the choices are made, the total adoption level is updated to $W_{t+1}$, which then becomes the starting point for the next round. This means your actions today don't change the state you are currently in, but they do set the stage for what happens next.

\subsubsection{Coordination Threshold and Stabilisation Dynamics}

The parameter $\theta \in (0,1)$ represents the level of cumulative adoption at which climate stabilisation benefits are strongest. Rather than modelling coordination success as a binary event, the model uses a smooth sigmoid function centred at $\theta$:
\[
  S(W) = \frac{1}{1 + \exp\!\bigl(-\eta(W - \theta)\bigr)}
\]
where $\eta > 0$ controls the steepness of the transition. This function takes values near zero when $W$ is well below $\theta$, equals 0.5 at $W = \theta$, and approaches 1 when $W$ is well above $\theta$.

The sigmoid captures a core feature of climate science: stabilisation benefits are not all-or-nothing. Partial coordination delivers partial benefit, but the marginal returns to coordination are highest around the threshold --- each additional adopter matters most when the world is near the tipping point.

The steepness parameter $\eta$ determines how concentrated benefits are around $\theta$:
\begin{itemize}
  \item High $\eta$: narrow transition window where most benefit is realised (tipping point dynamics).
  \item Low $\eta$: benefits spread more gradually across the full range of adoption.
  \item $\eta \to \infty$: recovers the hard indicator $\mathbf{1}\{W \geq \theta\}$ as a special case.
\end{itemize}

The value of $\theta$ itself is a modelling choice. A higher $\theta$ represents a world where broad coalition participation is necessary; a lower $\theta$ represents one where a smaller coalition of major players can stabilise the climate. Both $\theta$ and $\eta$ can be varied for comparative statics.

\subsection{Flow Payoffs and Their Components}

This model aims to capture some of the core costs and benefits of adopting or delaying adoption, and how these change over time alongside the benefits of successful cooperation. These are the components of the flow payoffs:

\begin{itemize}
  \item \textbf{Costs} --- How much does adoption cost a player today? Should I buy in now or wait for costs to fall as more players adopt?
  \item \textbf{Climate Damages} --- Whilst coordination is insufficient, countries pay for climate damages: increased extreme weather events, ecological damages, etc. These diminish smoothly as adoption approaches the threshold.
  \item \textbf{Political Pressure} --- Coordinating nations exert political pressure and impose costs on non-adopters. The strength of this pressure grows with coalition size.
  \item \textbf{Stabilisation Benefit} --- As coordination approaches and crosses the threshold, players increasingly enjoy the benefits of climate stabilisation. Damages diminish and the world moves away from catastrophic, irreversible harm.
\end{itemize}

For each of these components, the following sections break down how they are constructed and represented in the model, and how they create the full flow payoffs.

\subsubsection{Costs}

\textit{Depends on $W$ only.}
\[
  c_i(W) = c_i^0\,(1 - \gamma W)
\]
Adoption involves a one-off fixed cost. Costs fall with adoption. Players are assigned a cost factor $c_i^0$ which determines the baseline cost of renewable adoption; the more expensive initial adoption is, the higher $c_i^0$.

$\gamma \in [0,1]$ modulates how fast costs fall for all players with cumulative adoption, capturing learning-by-doing, supply chain development, and technology spillovers. Early adopters pay more; late movers benefit from the groundwork laid by others.

\subsubsection{Climate Damages}

\textit{Depends on $W$ and time.}
\[
  d_i(t, W) = d_i^0 \cdot (1 + \kappa t) \cdot \bigl(1 - S(W)\bigr)
\]
Countries face climate damages that diminish smoothly as cumulative adoption approaches and passes the coordination threshold. The damage function has three components:
\begin{itemize}
  \item \textbf{Baseline vulnerability} $d_i^0$: Each player's inherent exposure to climate harm. Economies dependent on agriculture, water resources, and outdoor activity face higher baseline damages.
  \item \textbf{Temporal escalation} $(1 + \kappa t)$: Damages grow each period the world remains insufficiently coordinated. The parameter $\kappa > 0$ introduces a ``deadline effect'' --- the longer coordination takes, the worse the consequences. This captures the reality that cumulative emissions and ecological degradation compound over time.
  \item \textbf{Coordination relief} $(1 - S(W))$: As adoption grows, the sigmoid $S(W)$ rises from near zero toward one, and damages correspondingly fall. When $W$ is far below $\theta$, $S(W) \approx 0$ and damages are at full intensity. When $W$ passes $\theta$, $S(W) \approx 1$ and damages are nearly eliminated. The transition is smooth, reflecting the fact that partial coordination delivers partial relief.
\end{itemize}

Crucially, damages are \textbf{non-excludable}. A player cannot escape them through individual adoption; only the collective level of adoption $W$ matters. This is what makes the model a cooperation problem rather than an individual optimisation.

\subsubsection{Political Pressure}

\textit{Depends on $W$ only.}
\[
  p_i \cdot W \cdot \mathbf{1}\{G_{it} = 0\}
\]
Non-adopters face political pressure that grows with the size of the adopting coalition. Each player has a pressure sensitivity $p_i > 0$ reflecting exposure to trade restrictions, diplomatic isolation, reputational costs, and domestic political backlash.

As more players adopt, holding out becomes increasingly costly. A country that refuses to join a coalition representing 10\% of global adoption faces mild pressure; one that holds out against 70\% faces intense isolation. This creates a bandwagon dynamic --- once adoption reaches a critical mass, the pressure on remaining holdouts accelerates.

\subsubsection{Stabilisation Benefit}

\textit{Depends on $W$ only.}
\[
  b \cdot S(W)
\]
As cumulative adoption approaches and crosses the threshold, all players receive an increasing per-period benefit. $S(W)$ rises smoothly from near zero (when adoption is low) to near $b$ (when adoption is high), with the steepest increase around $\theta$.

This captures the payoff from successful coordination: climate damages cease and the world avoids catastrophic, irreversible harm. The benefit is \textbf{non-excludable} --- even players who never adopted enjoy it as the world stabilises. This preserves the free-rider problem at the heart of climate cooperation.

Unlike a hard threshold where benefit switches on instantaneously, the sigmoid formulation reflects the reality that climate stabilisation is a continuous process. Some benefits accrue from partial coordination (e.g., reduced rate of warming, fewer extreme weather events), but the strongest benefits --- avoiding tipping points such as permafrost collapse and ice sheet destabilisation --- are concentrated around the threshold region.

\subsubsection{How Do These Combine?}

\noindent\textbf{Flow payoff from the current state} (everyone receives this regardless of action):
\[
  r_i(t, W) = -d_i(t, W) + b \cdot S(W)
\]
All players suffer climate damages and enjoy stabilisation benefits based on the state of the world. Both components now vary smoothly with $W$ through the sigmoid.

\noindent\textbf{Flow payoff from adopting}:
\[
  r_i^A = -c_i(W) + r_i(t, W)
\]
Adopters pay adoption costs and face climate damages/benefits based on the state.

\noindent\textbf{Flow payoff from waiting}:
\[
  r_i^D = -p_i W + r_i(t, W)
\]
Non-adopters do not pay adoption costs but face political pressure and climate damages/benefits based on the state.

\subsection{Recursive Value Structure}

The recursive value structure equations are broken down into two components: current flow payoffs and discounted future payoffs. These are the Q-values for adopting and waiting, and they dictate how players make decisions at any given state.

The value of adopting now:
\[
  Q_i^A(s) = r_i^A + \delta_i\,\mathbb{E}\bigl[V_i(s') \mid a_i = 1\bigr]
\]
The value of delaying:
\[
  Q_i^D(s) = r_i^D + \delta_i\,\mathbb{E}\bigl[V_i(s') \mid a_i = 0\bigr]
\]
Flow payoffs plus discounted value of future states, in short. Discounted values of future states are obtained via a recursion; this is possible due to the game being finite in nature.

\noindent\textbf{Terminal Condition:}
\[
  V_i(T+1,\,\cdot) = 0
\]

\noindent\textbf{Continuation Values (``Value of a State''):}

For players still deciding:
\[
  V_i(s) = \frac{1}{\lambda}\log\!\left(e^{\lambda Q_i^A} + e^{\lambda Q_i^D}\right)
\]
For players who have already adopted:
\[
  V_i(s) = r_i(W, t) + \delta_i\,\mathbb{E}[V_i(s')]
\]

Given the Q-values for each action, we need to represent the value of being at a particular state before a choice is made. This is the inclusive value. The inclusive value is what gets stored in the lookup table. When earlier periods compute expected continuation values, they reference these stored values, each summarising the entire future from that state onwards.

The formula computes a soft maximum over the two options, blending both in a way that favours the better one. This reflects how QRE players behave and captures optionality: the value of being at a state is higher than the weighted average of outcomes because you get to observe and then decide.

\subsection{Equilibrium}

\noindent\textbf{QRE:}
\[
  \sigma_i(s) = \frac{\exp\{\lambda Q_i^A\}}{\exp\{\lambda Q_i^A\} + \exp\{\lambda Q_i^D\}}
\]

Equilibrium in the model is based on Quantal Response Equilibrium. This avoids the knife-edge solutions of Nash equilibrium, where tiny parameter changes can flip outcomes entirely, and ensures clean comparative statics.

The formula converts Q-values into adoption probabilities. A higher value of adopting relative to delaying means a higher probability of adoption. The parameter $\lambda$ controls how sharply players respond: high $\lambda$ means near-perfect rationality, low $\lambda$ means more randomness.

At every state in the recursion, we solve for fixed-point strategies. This means we iterate until beliefs match behaviour: each player computes their best response given what they believe others will do, then we check if those beliefs were correct. If not, we update beliefs and repeat. When the probabilities players assume equal the probabilities that come out of everyone's optimisation, we have equilibrium. These equilibrium probabilities are stored in the lookup table and used to compute expected continuation values for earlier periods.

\noindent\textbf{Fixed Point:}
\[
  \sigma^* = \operatorname{QRE}\!\bigl(Q(\sigma^*)\bigr)
\]

An equilibrium is a set of adoption probabilities that are self-confirming. Each player's strategy is optimal given others' strategies, and others' strategies are optimal given theirs. The fixed point is where this circle closes: the probabilities we assume when computing Q-values are the same probabilities that come out of the QRE formula. No player wants to revise their beliefs or behaviour. We find this by iterating at each state until the probabilities stabilise.

\subsection{State Transitions}

After players make their choices, the state updates. Each player who adopts adds their weight to cumulative adoption. The new state becomes the starting point for the next period.

From any player's perspective, the next state is uncertain because others choose simultaneously. Each other player independently adopts or delays with some probability. The chance of landing in any particular future state is just the product of these independent probabilities. If you think there's a 60\% chance the US adopts and a 40\% chance the EU adopts, the chance both adopt is $0.6 \times 0.4$.

\subsection{Recursion Walkthrough}

\subsubsection*{Period $T$ (Terminal Period)}
\begin{itemize}
  \item No future to discount, so Q-values equal flow payoffs.
  \item For each of the possible configurations, compute flow payoffs for all players.
  \item Apply QRE formula directly (no iteration needed --- Q-values don't depend on others' strategies as there is no future round).
  \item Compute continuation values using the log-sum formula for active players.
  \item Store equilibrium probabilities $\sigma_i(T, G)$ and continuation values $V_i(T, G)$ in the lookup table.
\end{itemize}

\subsubsection*{Period $T - 1$}
\begin{itemize}
  \item Q-values now have two components: current flow payoffs $+$ discounted expected future value.
  \item Flow payoffs: computed for each of the 16 configurations.
  \item Expected future value: requires knowing (1) values of period $T$ states [stored], and (2) probabilities of reaching those states [unknown].
\end{itemize}

Finding equilibrium at each $T-1$ state through fixed-point iteration:
\begin{enumerate}
  \item Start with initial guess about what all active players will do (e.g., $\sigma = 0.5$ for all).
  \item Each player computes their $Q^A$ and $Q^D$ given these guesses.
  \item Apply QRE formula to each player's Q-values $\to$ get new probabilities.
  \item If new probabilities match the guesses $\to$ converged (fixed point found).
  \item If not $\to$ update guesses to the new probabilities and repeat.
\end{enumerate}

Once converged: equilibrium probabilities $\sigma^*$ are self-confirming (beliefs $=$ behaviour). Use $\sigma^*$ to compute transition probabilities via independent Bernoulli trials. Compute continuation values $V$ using final Q-values. Store $\sigma_i(T-1, G)$ and $V_i(T-1, G)$ for all configurations.

\subsubsection*{Period $T - 2$}

Same process as $T-1$, but now we look up stored $T-1$ values. Iterate to find equilibrium at each $T-2$ configuration. Store results and continue backward.

\subsubsection*{Why Recursion? Why QRE?}

\noindent\textbf{The temporal problem: optimising over time.}
Players don't simply maximise current payoffs --- they optimise their entire future path. A country deciding whether to adopt in period 5 must account for how that choice affects its options and payoffs in periods $6, 7, \ldots, T$. This creates a dynamic optimisation problem where current decisions have compounding consequences.

Recursion solves this by working backwards from the final period. At period $T$, there is no future, so the decision is simple. At period $T-1$, we can evaluate choices by looking one period ahead to $T$, whose values we've already solved. At $T-2$, we look ahead to $T-1$, which already incorporates information about $T$. By recursively embedding future values into current decisions, we solve the full dynamic problem without explicitly tracking every possible path through the game tree.

\noindent\textbf{The strategic problem: circular dependencies.}
My best strategy depends on your strategy, but your best strategy depends on mine. In climate cooperation, this circularity is acute: I should adopt early if I expect others to follow (costs fall, threshold reached). But I should wait if I expect others to free-ride (I pay costs, threshold missed). Everyone faces this same reasoning, creating a coordination problem with no obvious solution.

QRE solves this through fixed-point iteration. We guess what everyone will do, compute each player's best response to that guess, and check if the responses match the guess. If not, we update the guess and repeat. When beliefs finally equal behaviour --- when what I expect you to do is exactly what you want to do given what you expect from me --- we've found equilibrium. The probabilities are self-confirming.

QRE also has the added benefit of avoiding multiple equilibria that are knife-edge, allowing for cleaner comparative statics and parameter experimentation.

% ============================================================
\section{Calibration}
% ============================================================

Implementing the model requires assigning numerical values to all parameters. Some components of this parameterization can be disciplined directly by observable data, while others cannot. This identification challenge is common in game theoretic models of international cooperation, and the calibration strategy addresses it by separating the problem into two distinct layers.

The first layer is ordinal structure, which captures the relative positioning of blocs within each parameter channel. Relationships such as which bloc faces the highest transition costs or which is most vulnerable to climate damages can be inferred directly from data. World Bank development indicators, national emissions inventories, and trade statistics provide consistent evidence on the relative positions of the United States, European Union, China, and the Rest of the World. These cross bloc rankings are therefore treated as empirically identified and held fixed throughout the analysis.

The second layer is cardinal scaling, which determines the absolute magnitude of each payoff channel and their relative importance in shaping strategic behavior. The weight governments place on transition costs relative to climate damages is not directly observable and cannot be recovered from macro data alone. Rather than treating these channel weights as free scenario assumptions, they are estimated using a simulated method of moments (SMM) procedure.

Specifically, the vector of cardinal scaling parameters is chosen so that the model reproduces four stylized moments from the contemporary climate transition. These include the observed renewable adoption ratios between the EU and its peers, the near certain probability of the EU acting as a first mover, and the absence of near term global coordination cascades. By matching the model to these empirical moments, the calibration recovers the latent relative strength of costs, damages, and political pressure consistent with observed strategic behavior.

Accordingly, ordinal structure is data derived while cardinal scaling is structurally estimated. Subsequent sensitivity analysis is conducted around this fitted baseline, allowing the results section to distinguish local comparative statics from deeper structural mechanisms.

\subsection{Layer 1: Ordinal Structure}

\subsubsection{Per-Player Parameters}

Each per-player parameter is constructed from bloc-level data through a structural transformation anchored to a reference bloc (the EU). This approach preserves empirical ordinal relationships across blocs while allowing a single scaling parameter (alpha) to determine the absolute level.

Transition costs are derived from the carbon intensity of GDP, adjusted for development. Carbon intensity captures the extent of structural change required for decarbonisation, as higher-intensity economies must replace more emissions-intensive capital. However, low carbon intensity in less developed economies often reflects limited industrialisation rather than genuinely clean infrastructure. To account for this, the model applies a development-gated interpolation; a bloc's effective carbon intensity is defined as a weighted average of its observed intensity and a global baseline, where the weight depends on GDP per capita relative to the richest bloc. Developed economies retain their observed advantage, while less developed economies are partially adjusted toward the baseline. An additional affordability penalty increases costs for poorer blocs, reflecting higher financing constraints and weaker deployment capacity (Ameli et al., 2021).

\textbf{Climate damages:} are proxied by the share of agriculture, forestry, and fishing in GDP, capturing structural economic exposure to climate risk. Economies with larger primary sectors face greater direct losses from temperature extremes, water stress, and ecosystem disruption. While more composite vulnerability indices exist, agriculture share is preferred for its direct observability and consistency across blocs in World Bank national accounts data.

\textbf{Political pressure:} is proxied by trade openness, measured as trade as a percentage of GDP. This captures the degree to which a bloc's economic welfare depends on maintaining cooperative relations with major trading partners, and therefore its exposure to externally imposed costs from non-cooperation, including carbon border adjustments, climate-linked conditionality, and broader diplomatic friction.

\textbf{Influence weights:} are defined as an equal-weighted combination of emissions share and GDP share. Emissions share captures a bloc's necessity for any effective global agreement, while GDP share captures its economic leverage within the international system. Equal weighting reflects the assumption that neither dimension systematically dominates the other.

All parameters are normalised relative to the EU benchmark, so each $\alpha$ directly sets the EU's level for a given channel. The values for other blocs are determined by their data-implied ratios to the EU. This structure limits the modeller's degrees of freedom to the choice of $\alpha$, while the cross-bloc relationships remain fully data-driven.

\subsection{Layer 2: Global Parameters}

Global parameters govern the structure of the game rather than the characteristics of individual players. Unlike per-player parameters, they cannot be identified through cross-bloc comparisons. Instead, each is anchored to a baseline informed by published estimates in the climate economics and game theory literature.

These baselines serve as reference points rather than fixed values. Global parameters are treated as experimental inputs and varied systematically to examine how the underlying structure of the model shapes cooperation dynamics. In particular, they determine whether coordination emerges gradually or through sharp thresholds, how quickly costs decline through learning, and how strongly time preferences influence strategic behavior.

\subsubsection{Special Note on Player Discount Rates}

A limitation concerns time preferences. The social discount rate is one of the most consequential parameters in climate economics. Nordhaus (2007) adopts a descriptive approach based on observed market returns, implying discount rates of around 5.5\% per year and favouring gradual emissions reductions. Stern (2006), by contrast, takes a prescriptive stance grounded in intergenerational ethics, supporting rates closer to 1.45\% and advocating immediate, aggressive action. This is not a technical disagreement but a conceptual one; it turns on whether climate models should describe how governments behave or prescribe how they ought to behave. A game-theoretic framework concerned with strategic interaction aligns with the descriptive view. The relevant question is how blocs value the future in practice, not how they should. Accordingly, this framework adopts a descriptive stance.

Discount factors ($\delta$) are calibrated using a descriptive approach following Nordhaus (2007), reflecting observed political and economic time horizons rather than normatively optimal rates. Because each model period spans five years, $\delta$ represents the compounded effect of an annual discount rate over that interval. A discount value of $\delta = 0.75$, for example, corresponds to an annual rate of approximately 5.95\%, not 25\%.

Each bloc's discount factor combines pure time preference with institutional commitment credibility, interpreted as the probability that a climate commitment persists long enough to influence outcomes.

The European Union is assigned the highest value ($\delta = 0.85$, 3.25\% annually). The European climate law legally binds member states to net-zero by 2050 with intermediate enforcement mechanisms, while the EU ETS reflects over two decades of continuous policy. European Commission guidance recommending social discount rates of 3--4\% provides an additional institutional anchor.

China is assigned $\delta = 0.80$ (4.6\% annually). Centralised governance supports policy continuity through the five-year plan system and the formal embedding of ecological objectives. However, growth-related political constraints limit the effective horizon relative to the EU.

The United States receives $\delta = 0.75$ (5.9\% annually). The four-year electoral cycle introduces significant policy reversal risk, illustrated by withdrawal from and re-entry into the Paris agreement. The calibration therefore incorporates both standard time preference and a reversal premium.

The rest of the world is assigned $\delta = 0.65$ (9.0\% annually). Sovereign borrowing costs in many developing economies, particularly in sub-Saharan Africa, range between 6\%--10\%, and immediate fiscal and development pressures shorten effective planning horizons. This aligns with cost-of-capital estimates for clean energy investment in emerging markets (Ameli et al., 2021).

The ordering EU $>$ China $>$ US $>$ ROW is maintained across all specifications. The spread is chosen to generate meaningful strategic differentiation; the EU emerges as a natural early mover, China occupies an intermediate and potentially pivotal position, the US faces a sharper intertemporal trade-off, and the rest of the world remains the most constrained despite high vulnerability.

\subsubsection{Coordination Threshold}

The coordination threshold ($\theta = 0.80$) defines the level of weighted global participation at which stabilisation benefits become substantial. While the Paris agreement's 55\% emissions threshold governs entry into force, effective mitigation requires far broader participation. Net-zero pathways in IEA (2021) and IPCC 1.5\textdegree C scenarios consistently imply near-universal engagement among major emitters, with the four blocs modelled here collectively accounting for the substantial majority of global emissions. The chosen value reflects this near-universal requirement. As with other global parameters, $\theta$ is varied systematically to examine how the stringency of the coordination requirement affects equilibrium outcomes.

\subsubsection{Sigmoid Steepness}

The sigmoid steepness parameter ($\eta$) determines how sharply stabilisation benefits concentrate around the threshold. High values produce tipping-point dynamics, where benefits activate rapidly once $\theta$ is crossed. Lower values generate a more gradual accumulation of benefits across adoption levels. The baseline ($\eta = 15$) implies a steep transition, with $S(W)$ rising from approximately 0.06 to 0.94 within a $\pm 0.2$ band around $\theta$. Since this parameter governs the qualitative shape of payoffs rather than a directly observable quantity, it is treated as an experimental lever.

\subsubsection{QRE Rationality}

The quantal response equilibrium parameter ($\lambda = 1.5$) controls the sensitivity of behaviour to payoff differences. Experimental estimates typically range from 0.5 to 10 depending on context (McKelvey and Palfrey, 1995; Goeree, Holt, and Palfrey, 2003). The chosen value lies toward the lower end, reflecting the presence of noise in political decision-making arising from domestic constraints, incomplete information, and bureaucratic frictions. Higher values approach deterministic best-response behaviour; lower values increase randomness. Varying $\lambda$ allows the analysis to assess how behavioural noise influences coordination dynamics.

\subsubsection{Technology Learning Rate}

The learning rate ($\gamma = 0.25$) captures cost reductions driven by cumulative adoption through spillovers, supply chain development, and learning-by-doing. Empirical estimates suggest cost declines of 20--25\% per doubling of capacity for solar PV and 10--15\% for wind (Rubin et al., 2015; IRENA, 2023). The chosen value aligns with the upper end of these estimates, reflecting the fact that ``adoption'' in the model represents a broad technological transition rather than a single technology.

\subsubsection{Damage Escalation}

The damage escalation rate ($\kappa = 0.05$) governs how climate damages compound under delayed action. Over the model's 50-year horizon, this implies an increase of roughly 50\% in baseline damages under persistent non-cooperation. This magnitude is consistent with IPCC AR6 estimates suggesting losses on the order of 1--2\% of GDP per decade of delayed mitigation.

\subsubsection{Affordability Elasticity}

The affordability elasticity ($\varepsilon = 0.3$) determines how transition costs scale with income differences across blocs. Empirical evidence indicates that the cost of capital for renewable energy projects in developing economies can be two to three times higher than in advanced economies (Ameli et al., 2021). The chosen value reproduces cost differentials of a similar order of magnitude within the model.

\subsection{Layer 2: Simulated Method of Moments (SMM) Calibration}

The cardinal scaling layer of the model is estimated using a simulated method of moments procedure. The objective is to recover the latent relative strength of the four core payoff channels: cost scaling, damage scaling, pressure scaling, and benefit scaling. These represent the principal remaining degrees of freedom after the ordinal cross-bloc structure has been fixed by data transformations.

All other structural parameters are held fixed during estimation. This includes the bloc-specific discount factors, the rationality parameter, the sigmoid steepness, and the damage escalation. Holding these parameters fixed ensures that the estimated cardinal weights reflect the relative strength of strategic channels conditional on a common institutional environment rather than absorbing variation from unrelated structural assumptions.

For any candidate parameter vector, the dynamic game is solved under quantal response equilibrium using backward induction. The fitted equilibrium strategies are then passed through Monte Carlo simulation to generate model-implied moments. The SMM estimator selects the parameter vector that minimises the weighted squared distance between these simulated moments and their empirical targets.

\subsubsection{Targeted Moments}

The estimation targets four moments chosen to capture the most policy-relevant features of contemporary climate coordination.

\begin{itemize}
  \item \textbf{Moment 1: The EU--US Adoption Gap.} This targets the difference in period-one adoption probabilities between the European Union and the United States. It is designed to discipline the balance between transition costs and strategic patience. The empirical motivation is that the EU has consistently moved more aggressively than the US despite both possessing similar economic weight.

  \item \textbf{Moment 2: The Strategic Free-Rider Moment.} This captures the probability that Europe acts as an early mover while China delays and benefits from downstream spillovers. This moment is particularly informative for identifying the pressure and benefit channels because it disciplines how strongly early leadership translates into leverage over pivotal holdouts.

  \item \textbf{Moment 3: The EU Leadership Ratio.} This measures how much more likely the EU is to adopt in the first period relative to its strongest rival. It identifies the degree to which the fitted economy reproduces the role of Europe as the natural first mover in global climate diplomacy.

  \item \textbf{Moment 4: Global Threshold Success.} This moment targets strategic delay. Empirically, global climate coordination does not immediately cascade even when early movers adopt. The target is set at a near-zero value of 0.02 to prevent the objective surface from becoming flat while still enforcing delayed tipping as the dominant equilibrium pattern.
\end{itemize}

\begin{figure}[htbp]
  \centering
  \includegraphics[width=0.75\textwidth]{results/dissertation/table_smm_alphas.png}
  \caption*{Table: Calibrated Cardinal Scaling Parameters ($\hat{\alpha}$). Values recovered by the SMM estimator; all structural and ordinal parameters held fixed during estimation.}
  \label{tab:smm_alphas}
\end{figure}

Because four moments are used to identify four unknown cardinal weights, the system is exactly identified. The objective is not to overfit a large moment set but to recover the single parameter vector that best reproduces the key stylized facts the model is designed to explain.

\subsubsection{Verification and Stability}

The optimisation is run from multiple starting points to serve as a convergence diagnostic. If all starting points converge to the same solution, the fitted baseline is interpreted as a unique local minimum. Following estimation, three verification checks are applied.

First, optimiser stability is tested by restarting the numerical search from the estimated optimum to confirm the solution remains unchanged. Second, Monte Carlo seed stability is assessed by re-running the fitted model across multiple random seeds at high simulation depth. The fitted economy is only carried forward if coordination success rates remain tightly concentrated, ensuring that results are not driven by simulation noise. Finally, ordinal ranking preservation verifies that estimation does not overturn the data-implied cross-bloc orderings. The fitted economy must preserve the empirical rankings for costs, damages, and political pressure. This is crucial because the core contribution of this dissertation rests on combining data-disciplined ordinal structure with structurally estimated cardinal channel strength.

% ============================================================
\section{Experimental Design}
% ============================================================

\subsection{What This Model Is For}

The model is not designed as a forecasting tool. It does not attempt to predict whether real-world climate coordination will succeed or fail, and its outputs should not be interpreted as literal probabilities of observed outcomes. Instead, its purpose is mechanism exploration. Starting from the SMM-fitted strategic environment, the objective is to identify which policy channels exert the greatest leverage over coordination dynamics and whether those rankings remain stable once multiple structural assumptions are perturbed.

The experimental design follows directly from this objective. Each experiment begins from the fitted baseline economy, perturbs one or more structural parameters, and records how coordination success, timing, and bloc-level behavior shift in response. The shape of these responses reveals the relative sensitivity of the system to different strategic levers and provides a policy map of which channels can materially reduce coordination risk.

\subsection{Marginal Parameter Sweeps}

The primary analytical tool is the continuous marginal parameter sweep. One parameter is varied continuously across a predefined range while all remaining parameters are held at their SMM-estimated baseline values. At each point, the full dynamic game is re-solved and Monte Carlo simulation records the coordination success rate. The result is a marginal success curve with the parameter value on the horizontal axis and the probability of successful coordination on the vertical axis.

These curves answer a fundamental policy question regarding how far a given strategic lever must move before the system changes regime. A steep curve implies that small deviations from the fitted baseline generate large changes in coordination probability. Conversely, a flat curve implies that the channel has limited leverage even under substantial perturbation.

Two summary statistics characterise each curve. The 90\% crossing is the percentage deviation from the SMM-fitted baseline required to raise coordination success above 90\%, while the 95\% crossing is the deviation required to make coordination failure very unlikely. These thresholds are chosen to reflect the tail-risk framing of the analysis: the relevant policy question is not merely whether coordination becomes more likely than not, but whether the residual probability of catastrophic climate delay is compressed to a small tail. The distance from the baseline to each crossing provides a direct measure of policy leverage. If the 90\% crossing lies close to zero, the channel is highly sensitive and relatively small interventions are sufficient to improve outcomes. If the crossing is distant or absent, the channel has weak leverage. When a crossing is reported as absent, it means that no plausible deviation of that parameter is sufficient to move the fitted economy past the relevant policy frontier.

The channels swept include both bloc-specific and global structural levers: discount factors (individually and jointly for the United States and China), cost multipliers, pressure multipliers, the technology learning rate, and the coordination threshold. Together, these sweeps provide a detailed elasticity map of the fitted climate coordination economy and identify which strategic levers are capable of shifting the system toward safer equilibrium regions.

\subsection{Global Sensitivity Analysis}

Marginal sweeps isolate one channel at a time. While this cleanly identifies local leverage, it cannot establish whether these rankings remain robust once multiple structural forces move simultaneously. The global sensitivity analysis addresses this by jointly perturbing all major structural parameters around the SMM-fitted baseline.

One thousand parameter vectors are drawn from uniform distributions spanning plausible ranges for the key structural inputs. These include the four bloc-specific discount factors, the technology learning rate, sigmoid steepness, QRE rationality, and damage escalation. For each draw, the full model is resolved and the key outcomes are recorded, including coordination success and mean coordination timing.

Spearman rank correlations between each input parameter and each outcome provide a non-parametric measure of global influence across the full parameter space. A parameter that remains highly ranked in both marginal sweeps and the global sensitivity analysis can be interpreted as a \textbf{structurally robust policy lever}. By contrast, a parameter that appears powerful in isolation but weak under joint variation suggests strong interaction effects or local nonlinearity around the fitted baseline.

The purpose of the global sensitivity analysis is therefore robustness validation rather than scenario comparison. It tests whether the local policy rankings implied by marginal sweeps survive once the wider strategic environment is allowed to move, thereby separating genuinely robust findings from effects that depend on holding the rest of the system fixed.

% ============================================================
\section{Results}
% ============================================================

\subsection{Model Baselines}

Before examining parameter perturbations, the SMM-fitted baseline economy is first reported. The model retains the same ordinal cross-bloc structure derived from World Bank and related macroeconomic data, while the cardinal channel weights governing costs, damages, political pressure, and coordination benefits are now estimated through the simulated method of moments procedure. Under this fitted strategic environment, coordination succeeds in 62.8\% of simulated paths, with a mean coordination timing of 3.51 periods.

This gives concrete equilibrium form to the strategic intuition introduced at the outset. The fitted world is not one in which cooperation is impossible. Instead, coordination is more likely than failure, but the equilibrium remains fragile: a substantial residual probability of delayed or failed coordination persists. In other words, the blocs often do eventually move, but not always early enough, and not always at all. This baseline provides the interpretive anchor for all subsequent perturbation analysis. The purpose of the marginal sweeps and global sensitivity analysis is therefore to identify which strategic levers most effectively improve the fitted equilibrium along two margins: the timing margin, by bringing average coordination forward, and the risk margin, by compressing the remaining tail risk of catastrophic delay.

\begin{figure}[htbp]
  \centering
  \includegraphics[width=0.85\textwidth]{results/dissertation/fig_w_distribution.png}
  \caption{Distribution of Final Coalition Weight $W_T$ --- SMM Baseline ($n = 1000$ Monte Carlo runs). The threshold $\theta = 0.8$ is marked by the red dashed line; the mean $\bar{W} = 0.562$ by the orange dot-dash line.}
  \label{fig:w_dist}
\end{figure}

\begin{figure}[htbp]
  \centering
  \includegraphics[width=0.85\textwidth]{results/dissertation/table_baseline_diagnostics.png}
  \caption*{Table 1: SMM Baseline Monte Carlo Diagnostics ($n = 1000$, seed 42).}
  \label{tab:baseline}
\end{figure}

\begin{figure}[htbp]
  \centering
  \includegraphics[width=0.85\textwidth]{results/dissertation/fig_period1_sigmas.png}
  \caption{Period-1 Equilibrium Adoption Probabilities by Bloc --- SMM Baseline. The EU is the clear first-mover; China and RoW adoption probabilities are near zero in the initial period, consistent with the strategic delay targeted by Moment 4.}
  \label{fig:period1}
\end{figure}

\begin{figure}[htbp]
  \centering
  \includegraphics[width=\textwidth]{results/dissertation/fig_sigma_over_time.png}
  \caption{Evolution of Equilibrium Adoption Probabilities Across Periods --- SMM Baseline. EU leadership is concentrated in early periods; US and CN probabilities rise only as the coalition weight accumulates, reflecting the endogenous pressure mechanism.}
  \label{fig:sigma_time}
\end{figure}

\begin{figure}[htbp]
  \centering
  \includegraphics[width=\textwidth]{results/dissertation/fig_adoption_timing.png}
  \caption{Distribution of Adoption Timing by Bloc --- SMM Baseline ($n = 1000$ Monte Carlo runs). The EU adopts predominantly in early periods; the US and CN timing distributions are right-skewed, with substantial mass on never-adopting paths, reflecting the residual failure probability.}
  \label{fig:adopt_timing}
\end{figure}

\FloatBarrier

\subsection{Marginal Sweeps}

The marginal sweep analysis varies one parameter at a time across a continuous range, holding all others at baseline, and records the coordination success rate at each point. The 80\% crossing, reported as the percentage deviation from the calibrated baseline required to raise coordination success above 80\%, and the 95\% crossing, the deviation required to make coordination near-certain, summarise each channel's leverage.

\begin{figure}[htbp]
  \centering
  \includegraphics[width=\textwidth]{results/dissertation/fig_key_sweeps.png}
  \caption{Marginal Sensitivity Sweeps (Key Channels) --- SMM Baseline. Each panel varies one structural parameter continuously while holding all others fixed. Red dashed line marks the 80\% coordination frontier; blue dotted line marks the 95\% frontier. Shaded region indicates the zone of near-certain coordination.}
  \label{fig:key_sweeps}
\end{figure}

\begin{figure}[htbp]
  \centering
  \includegraphics[width=\textwidth]{results/dissertation/table_sweep_crossings.png}
  \caption*{Table 2: Marginal Sweep Threshold Crossings (SMM Baseline). Crossing values expressed as \% change from the SMM-calibrated baseline. Positive values indicate the parameter must increase from baseline; negative values indicate it must decrease. --- indicates crossing absent within sweep range.}
\end{figure}

\subsubsection{Marginal Sweep Findings}

The marginal sweep frontiers reveal a clear hierarchy of strategic leverage around the SMM-fitted baseline, separating levers that make coordination reliably likely from those capable of pushing failure into a remote tail event.

The strongest channels are those linked to strategic patience and coalition stringency. The coordination threshold is the single most sensitive lever, requiring only a modest reduction to move the fitted economy above the 90\% frontier and a further small reduction to push success above 95\%. Among bloc-specific parameters, the most powerful channels are the United States discount factor, the China discount factor, and the joint US--China patience channel. The fact that the China and joint patience channels cross both frontiers at the same point is especially revealing: once these pivotal blocs become sufficiently forward-looking, the system does not merely become safer, but rapidly transitions into a near-certain coordination regime. The fact that the China and joint patience channels cross both frontiers at the same point is especially revealing: once these pivotal blocs become sufficiently forward-looking, the system does not merely become safer, but rapidly transitions into a near-certain coordination regime. This indicates that the fitted equilibrium sits closest to a tipping boundary along the dimensions of pivotal-bloc patience and effective coalition requirements.

Transition costs form the second tier of leverage. A modest 6.36\% reduction in the global cost multiplier is sufficient to move the system above 80\%, while a 15.45\% reduction pushes it above 95\%. Bloc-specific cost reductions are also meaningful, though somewhat weaker. China's cost multiplier requires a 15.45\% reduction for the 80\% frontier and 26.36\% for 95\%, while the United States requires 15.45\% and 37.27\% respectively. The Rest of the World remains the least efficient cost lever, requiring much larger reductions. These results suggest that cost declines matter most as secondary accelerants, improving reliability once the patience and threshold conditions are already close to tipping.

The slower-moving structural feedback channels are substantially weaker. Learning spillovers require an 85.44\% increase to cross the 80\% frontier and 156.36\% for 95\%, while damage accumulation requires 92.80\% and a near-extreme 300\% increase respectively. These channels clearly strengthen coordination dynamics, but their distance from both frontiers suggests they operate through long-run amplification rather than immediate control over baseline tail risk.

The most striking finding concerns pressure-based channels. Neither the US nor China pressure multipliers generate an 80\% or 95\% crossing anywhere within the tested range, while the global pressure multiplier requires an extreme 196.36\% increase merely to cross the lower frontier and never reaches 95\%. This indicates that pressure is not an efficient standalone intervention when varied in isolation. Instead, its role appears fundamentally interaction-dependent: pressure amplifies coordination once other conditions begin to shift, but does not by itself move the fitted equilibrium into a safe region.

Overall, the sweep results point to a sharp policy hierarchy. The most effective way to improve the fitted equilibrium is through modest increases in the forward-looking credibility of pivotal blocs, particularly the United States and China jointly, combined with small reductions in effective coalition stringency. Cost reductions provide strong secondary support, while pressure and long-run feedback channels appear to matter primarily as complementary amplifiers rather than first-order standalone levers.

\subsection{Global Sensitivity Analysis}

The global sensitivity analysis (GSA) estimates Spearman rank correlations between the full set of major structural input parameters and the key outcome dimensions of the fitted economy: coordination success rate, coordination timing, and first-mover probabilities. These correlations are computed across one thousand joint parameter draws centred on the SMM-fitted baseline, with bloc-specific costs, discount factors, threshold structure, behavioural parameters, and pressure channels all allowed to vary simultaneously within plausible ranges. The purpose of the exercise is to test whether the local leverage rankings identified by the marginal sweeps remain robust once the broader strategic design space is allowed to move.

\begin{figure}[htbp]
  \centering
  \includegraphics[width=\textwidth]{results/dissertation/fig_gsa_sensitivity.png}
  \caption{Global Sensitivity Analysis --- Spearman Rank Correlations. Each bar shows the correlation between a structural input parameter and coordination success (left panel) and mean coordination timing (right panel) across 1000 joint parameter draws. Parameters are ordered by absolute correlation with coordination success.}
  \label{fig:gsa}
\end{figure}

\begin{figure}[htbp]
  \centering
  \includegraphics[width=\textwidth]{results/dissertation/fig_gsa_heatmap.png}
  \caption{Global Sensitivity Analysis --- Heatmap of Spearman $\rho$ across all input--output pairs. Rows are structural parameters; columns are outcome dimensions (coordination success, mean timing, first-mover probabilities). Colour intensity indicates the magnitude of the rank correlation; sign is indicated by hue.}
  \label{fig:gsa_heat}
\end{figure}

\begin{figure}[htbp]
  \centering
  \includegraphics[width=0.9\textwidth]{results/dissertation/table_gsa_correlations.png}
  \caption*{Table: GSA Spearman Rank Correlations --- Full Results. Correlations with $|{\rho}| > 0.1$ marked; $p < 0.001$ for all entries exceeding this threshold given $n = 1000$ draws.}
  \label{tab:gsa_corr}
\end{figure}

\subsubsection{GSA Findings}

GSA reveals a richer hierarchy than the local sweeps alone, showing that the fitted coordination economy is jointly governed by coalition threshold stringency, China's domestic transition burden, and the patience of the major emitters.

The single strongest negative correlate of coordination success is the coalition threshold $\theta$ (rho = $-0.415$, $p < 0.001$). Higher participation requirements sharply reduce coordination reliability and materially delay tipping once blocs begin to move. This strongly reinforces the marginal sweep finding that threshold geometry is one of the most efficient strategic levers in the fitted economy. It also strengthens its substantive interpretation: institutions that allow partial stabilisation gains to become salient before near-universal participation can substantially compress residual failure risk.

The second dominant block of leverage comes from transition costs, especially China's cost channel. Cost\textsubscript{CN} is strongly negatively associated with coordination success (rho = $-0.366$, $p < 0.001$) and exhibits the largest negative correlation with China's first-mover probability (rho = $-0.644$, $p < 0.001$). This is highly revealing. Once bloc-specific costs are allowed to vary directly, the remaining coordination risk is no longer purely a patience problem, but a China-specific cost-and-credibility bottleneck. China's willingness to move first is jointly constrained by its effective planning horizon and its domestic transition burden.

The patience channels remain strongly positive and globally robust. $\delta_{CN}$ (rho = $+0.271$, $p < 0.001$) and $\delta_{US}$ (rho = $+0.231$, $p < 0.001$) continue to rank among the strongest positive correlates of coordination success, while both materially accelerate coordination timing. This confirms that the central sweep result survives the expanded parameter space: the US--China credibility axis remains a first-order driver of whether blocs coordinate early enough to avoid the failure tail.

A particularly important result is the near-total irrelevance of explicit pressure channels. Pressure\textsubscript{CN}, Pressure\textsubscript{US}, Pressure\textsubscript{EU}, and Pressure\textsubscript{RoW} all exhibit correlations close to zero across success, timing, and first-mover outcomes. This is now a much stronger negative finding than in the narrower GSA because these channels are explicitly allowed to vary over broad ranges. The implication is not that pressure is unimportant to the model's internal dynamics, but that pressure works primarily as an endogenous emergent consequence of movement in other channels rather than as an independently tunable standalone lever.

The behavioural precision parameter $\lambda$ again primarily shapes the path rather than the endpoint. Its effect on coordination success remains weak, but it strongly accelerates coordination timing (rho = $-0.347$, $p < 0.001$). This suggests that greater strategic precision sharpens the tipping cascade once blocs become unstable, even if it does little to determine whether instability emerges in the first place.

Overall, the expanded GSA sharpens the interpretation of the fitted baseline. The residual coordination risk is best understood as a threshold-and-China bottleneck problem, in which effective coalition stringency, China's domestic transition burden, and the long-horizon credibility of the major emitters jointly determine whether the system escapes catastrophic delay.

% ============================================================
\section{Policy Implications}
% ============================================================

The policy implications follow from two complementary layers of evidence. Marginal sweep frontiers identify the local parameter movements required to move the fitted baseline across the 90\% and 95\% coordination frontiers, while the expanded global sensitivity analysis tests whether these leverage rankings remain dominant once the wider strategic design space is jointly perturbed. Implications are stated most confidently where both local frontier distances and global multivariate importance point in the same direction.

\subsection{High-Confidence Core Levers}

The strongest and most robust implication concerns the planning horizons of the major emitters, particularly China and the United States. Both the marginal sweeps and the GSA identify $\delta_{CN}$ as the dominant coordination lever, with the United States consistently second. Small increases in Chinese patience are sufficient to move the fitted economy from a fragile 62.8\% success world into both the 80\% and 95\% frontiers, while the GSA confirms that $\delta_{CN}$ remains one of the strongest positive correlates of success and materially accelerates coordination timing. The coordination problem is therefore best understood not as a generic free-rider problem, but as a pivotal-holdout timing problem concentrated in the major emitters, especially China.

The joint US--China patience channel is especially powerful. The sweep frontiers show that simultaneous improvements in both blocs' planning horizons move the system into a near-certain coordination regime with only modest deviation from baseline. This strongly supports the policy value of bilateral credibility mechanisms, such as joint declarations, reciprocal ratchets, and institutional arrangements that reduce the probability of unilateral reversal. Bilateral climate diplomacy between the United States and China is therefore worth more than the sum of its parts because it directly targets the model's most sensitive strategic axis.

A second high-confidence implication concerns the perceived distance to successful cooperation. The strong influence of $\theta$ in both the sweeps and GSA suggests that what matters is not a literal threshold choice, but whether blocs perceive meaningful gains from joining partial coalitions before near-universal participation is reached. Policy architectures that make progress feel closer, through visible milestones, staged reward mechanisms, and early coalition wins, can therefore materially reduce the psychological and strategic distance to tipping.

\subsection{Second-Order Amplifiers}

The second class of levers consists of transition cost reductions, especially China's domestic transition burden. The expanded GSA shows that Cost\textsubscript{CN} is one of the strongest negative correlates of both success and China first-mover probability, indicating that the residual coordination risk is partly a China-specific cost bottleneck. This substantially strengthens the case for industrial policy, clean energy subsidies, grid investment, and financing structures that directly reduce the domestic burden of early movement for pivotal emitters.

More broadly, global and bloc-specific cost reductions produce relatively small 90\% and 95\% frontier distances, confirming that industrial policy remains an important tool for improving coordination reliability. However, both the frontier distances and the GSA rankings suggest that these channels are secondary accelerants rather than the primary trigger. Lower costs help blocs move once the strategic environment is already close to tipping, but do less to resolve the underlying holdout problem than improvements in long-horizon credibility.

The GSA further shows that sigmoid sharpness and damage escalation matter more globally than their marginal frontier distances alone suggest. This implies that urgency and the concentration of rewards around visible milestones mainly operate through interaction effects. Policies that make the costs of delay more immediate, whether through climate-risk disclosure, loss-and-damage accounting, or sectoral transition deadlines, are therefore best understood as amplifying the effectiveness of patience-targeting institutions rather than replacing them.

\subsection{Conditional Complements}

A third class of instruments appears powerful only in interaction with the broader strategic environment. The clearest example is political pressure. The marginal sweeps show that pressure-based channels are weak standalone levers, and the expanded GSA confirms near-zero correlations across success, timing, and first-mover outcomes. Yet pressure remains central to the model's endogenous mechanism because it compounds once other conditions begin to shift. This implies that tools such as carbon border adjustments, trade-linked climate standards, and diplomatic conditionality should be viewed as complements to a functioning credibility architecture rather than substitutes for it.

Technology learning spillovers follow a similar logic. Faster learning materially strengthens coordination cascades once adoption begins, but both the frontier distances and the GSA show that learning alone is too distant from the fitted baseline to serve as the primary risk-reduction instrument. It is therefore best interpreted as a cascade amplifier rather than the first-order solution to strategic delay.

\subsection{The Overarching Implication}

The broad policy lesson is that the path of least resistance is narrow but clear. The fitted strategic economy suggests that the highest-return interventions are those that extend the effective planning horizons of the United States and China, reduce China's domestic transition burden, make successful cooperation feel strategically closer, and use cost and pressure channels as reinforcing complements.

In substantive terms, the model implies that the central policy problem is not simply accelerating technology diffusion, but making pivotal blocs willing to move sooner and with greater confidence that progress is both meaningful and reciprocated. Once that axis shifts, the wider endogenous pressure and spillover mechanisms become capable of rapidly compressing the remaining tail risk of catastrophic delay.

% ============================================================
\section{Limitations}
% ============================================================

A number of structural limitations constrain the scope of the conclusions and point toward natural extensions.

\textbf{Homogeneous rationality.} The model assigns a single rationality parameter ($\lambda$) to all blocs, ignoring that the US political system, Chinese planning apparatus, and EU legislative process introduce fundamentally different frictions. Heterogeneous $\lambda$ values would likely disrupt the finding that strategic noise uniformly benefits coordination and could reveal asymmetric rationality as an independent source of delay.

\textbf{No belief updating.} Backward induction assumes blocs hold correct beliefs from the outset. Real climate diplomacy is largely a story of learning: the 2014 US--China joint announcement mattered because it was a public signal that shifted each party's beliefs about the other's type, not merely their discount factors. A Bayesian extension would distinguish agreements that shift preferences from those that resolve uncertainty about preferences, a distinction the current model erases.

\textbf{Sigmoid benefit function.} The threshold structure drives the model's tipping dynamics, but whether climate benefits actually concentrate near a participation threshold is an open empirical question. If benefits accrue more linearly, the structural advantage of continuous pressure over front-loaded rewards would attenuate, potentially reversing the key cross-scenario finding. The functional form is held fixed throughout, the assumption least grounded in direct evidence and most consequential for the results.

\textbf{History-independent payoffs.} The model treats each period as structurally independent. In practice, early adoption builds credibility that lowers future commitment costs; defection raises them. $\delta$ partially captures this, but conflates time preference with institutional credibility, which are distinct in a dynamic game. The patience-dominance finding is likely robust, but the model systematically understates the value of early signalling and commitment devices.

\end{document}
